\chapter{Harmonic Forms and the Hodge Theorem}\label{hodge-theorem}
\chapterstatus{complete}

\section{Introduction}\label{hodge-theorem:section-introduction}

On a compact oriented Riemannian manifold every de Rham cohomology class contains exactly one harmonic form
(Theorem~\ref{hodge-theorem:theorem-hodge}). This is the theorem that makes Hodge theory possible: it turns a topological object,
the cohomology of a manifold, into the solution space of an elliptic equation, on which the metric and, in the complex case, the
complex structure act. The idea is already visible in finite dimensions: for a complex of finite-dimensional inner product spaces, each cohomology class has a unique representative of
smallest norm, and these representatives are the kernel of a ``Laplacian'' $dd^* + d^*d$ (Lemma~\ref{hodge-theorem:lemma-finite-dimensional}). On a compact manifold the same statement holds for
the de Rham complex, and the only difficulty is analytic: one needs to know that the Laplacian has closed range and finite-dimensional kernel, which is what
Chapter~\ref{elliptic-operators} provides.

The consequences drawn here are the Hodge decomposition of the space of forms (with the Helmholtz decomposition of vector fields on a torus as a concrete instance), Poincar\'e duality via the
Hodge star, the conformal invariance of harmonic $1$-forms on surfaces, the Poincar\'e inequality, and Bochner's vanishing theorem, with a complete proof of the Bochner--Weitzenb\"ock formula on
$1$-forms; the complex analogue is in Chapter~\ref{kahler}. The later sections develop the theory further:
\begin{itemize}
\item Section~\ref{hodge-theorem:section-elliptic-complexes}: elliptic complexes, the Hodge theorem for them (finite-dimensional cohomology, harmonic representatives), twisted de Rham complexes, and the Euler
characteristic of a complex as an index;
\item Section~\ref{hodge-theorem:section-spectrum}: the positive eigenvalues of the Laplacian form exact complexes, the finite form of the McKean--Singer identity, the heat flow to the harmonic representative,
and Lichnerowicz's lower bound for the first eigenvalue under a Ricci curvature bound;
\item Section~\ref{hodge-theorem:section-middle-degree}: conformal invariance of the Hodge star in the middle degree, self-dual and anti-self-dual forms, and the signature as $b^+ - b^-$;
\item Section~\ref{hodge-theorem:section-products}: harmonic forms on Riemannian products.
\end{itemize} References: \cite[Chapter 6]{warner-foundations}, \cite[Chapter 4]{wells-differential-analysis},
\cite[Chapter 0]{griffiths-harris}, \cite[Chapter 5]{voisin-hodge-1}.

\noindent\textbf{Prerequisites.} Chapter~\ref{elliptic-operators} (Elliptic Operators on Compact Manifolds) and
Chapter~\ref{de-rham} (de Rham Cohomology).

Throughout, $(M, g)$ is a compact oriented Riemannian $n$-manifold, with the $L^2$ inner product, the codifferential $\delta$ and the
Laplacian $\Delta = d\delta + \delta d$ of Definition~\ref{riemannian-geometry:definition-laplacian}, and
$\mathcal{H}^k = \mathcal{H}^k(M, g) = \{\alpha \in \Omega^k(M) : \Delta\alpha = 0\}$.

\section{The Hodge decomposition}\label{hodge-theorem:section-decomposition}

We begin with the finite-dimensional model, where everything is linear algebra.

\begin{lemma}[Hodge theory in finite dimensions]\label{hodge-theorem:lemma-finite-dimensional}
Let $V^{k-1} \xrightarrow{d} V^k \xrightarrow{d} V^{k+1}$ be linear maps of finite-dimensional real inner product spaces with $d \circ d = 0$, let $\delta$ denote the adjoints of the two maps $d$, and
$\Delta = d\delta + \delta d \colon V^k \to V^k$. Then $\ker\Delta = \ker d \cap \ker\delta$, there is an orthogonal decomposition
\[
V^k = \ker\Delta \oplus dV^{k-1} \oplus \delta V^{k+1},
\]
and $\ker d = \ker\Delta \oplus dV^{k-1}$. Hence every class in $\ker d/dV^{k-1}$ has a unique representative in $\ker\Delta$, which is its unique element of smallest norm.
\end{lemma}

\begin{proof}
For $v \in V^k$, $\langle\Delta v, v\rangle = \langle\delta v, \delta v\rangle + \langle dv, dv\rangle$, so $\Delta v = 0$ implies $dv = 0$ and $\delta v = 0$; the converse is clear. The images of $d$ and $\delta$ are
orthogonal: $\langle du, \delta w\rangle = \langle ddu, w\rangle = 0$. In finite dimensions the orthogonal complement of an image is the kernel of the adjoint, so the orthogonal complement of
$dV^{k-1} \oplus \delta V^{k+1}$ is $\ker\delta \cap \ker d = \ker\Delta$, which gives the decomposition. If $dv = 0$ and $v = h + du + \delta w$, then $0 = dv = d\delta w$, so
$\|\delta w\|^2 = \langle w, d\delta w\rangle = 0$, and $v = h + du$. Finally $\|h + du\|^2 = \|h\|^2 + \|du\|^2$, which is smallest exactly for $du = 0$.
\end{proof}

For $\Omega^\bullet(M)$ the spaces are infinite-dimensional and two things could go wrong: the image of $d$ need not be closed (in $L^2$), so that ``orthogonal complement of the image = kernel of the
adjoint'' only holds up to closure, and the kernel of $\Delta$ could be infinite-dimensional. Ellipticity of $\Delta$ rules out both.

\begin{theorem}\label{hodge-theorem:theorem-decomposition}
For each $k$ the space $\mathcal{H}^k$ is finite-dimensional and consists of the forms with $d\alpha = 0$ and $\delta\alpha = 0$, and there is an orthogonal direct sum decomposition
\[
\Omega^k(M) = \mathcal{H}^k \oplus d\Omega^{k-1}(M) \oplus \delta\Omega^{k+1}(M),
\]
with $\ker d = \mathcal{H}^k \oplus d\Omega^{k-1}$ and $\ker\delta = \mathcal{H}^k \oplus \delta\Omega^{k+1}$. Moreover there is a unique operator $G \colon \Omega^k(M) \to \Omega^k(M)$, the \emph{Green operator},
that vanishes on $\mathcal{H}^k$, maps $(\mathcal{H}^k)^\perp$ into itself, and satisfies $\alpha = H\alpha + \Delta G\alpha$, where $H$ is the orthogonal projection onto $\mathcal{H}^k$; it commutes with $d$,
$\delta$ and $\Delta$.
\end{theorem}

\begin{proof}
$\Delta$ is elliptic, formally self-adjoint and non-negative (Example~\ref{elliptic-operators:example-elliptic} and
Proposition~\ref{riemannian-geometry:proposition-adjoint}), so by Theorem~\ref{elliptic-operators:theorem-fredholm} its kernel $\mathcal{H}^k$ is finite-dimensional and consists of smooth forms, and
\[
\Omega^k(M) = \mathcal{H}^k \oplus \Delta\Omega^k(M)
\]
orthogonally: indeed $L^2 = \operatorname{im}\Delta \oplus \ker\Delta$ by Theorem~\ref{elliptic-operators:theorem-fredholm}(3), and a smooth form orthogonal to $\mathcal{H}^k$ is $\Delta\beta$ for a smooth $\beta$
by elliptic regularity. The description of $\mathcal{H}^k$ as $\ker d \cap \ker\delta$ is
Proposition~\ref{riemannian-geometry:proposition-adjoint}(3).

Since $\Delta\beta = d(\delta\beta) + \delta(d\beta)$, the second summand is contained in $d\Omega^{k-1} + \delta\Omega^{k+1}$, so the three spaces span $\Omega^k$. They are pairwise orthogonal: harmonic forms are
orthogonal to $d\Omega^{k-1}$ because $(\alpha, d\beta) = (\delta\alpha, \beta) = 0$, and similarly to $\delta\Omega^{k+1}$; and $(d\beta, \delta\gamma) = (dd\beta, \gamma) = 0$. Hence the sum is direct and
orthogonal. If $d\alpha = 0$ and $\alpha = \eta + d\beta + \delta\gamma$, then $0 = (d\alpha, \gamma) = (d\delta\gamma, \gamma) = \|\delta\gamma\|^2$, so $\alpha \in \mathcal{H}^k \oplus d\Omega^{k-1}$; the statement for
$\ker\delta$ is symmetric.

For the Green operator, define $G\alpha = 0$ on $\mathcal{H}^k$ and, on the orthogonal complement, $G\alpha = \beta$ where $\beta$ is the unique solution of $\Delta\beta = \alpha$ orthogonal to $\mathcal{H}^k$
(Theorem~\ref{elliptic-operators:theorem-fredholm}(3) gives existence and smoothness, and uniqueness holds because two solutions differ by a harmonic form). Then $\alpha = H\alpha + \Delta G\alpha$ by
construction. Since $\Delta$ commutes with $d$, $\delta$ and $\Delta$, and these operators preserve $\mathcal{H}^\bullet$ and its orthogonal complement (for instance $d\mathcal{H}^k = 0$ and
$(d\alpha, \eta) = (\alpha, \delta\eta) = 0$ for $\eta$ harmonic), the uniqueness of $G$ forces $Gd = dG$, $G\delta = \delta G$ and $G\Delta = \Delta G$.
\end{proof}

Explicitly, the decomposition of a form $\alpha$ is
\[
\alpha = H\alpha + d(\delta G\alpha) + \delta(dG\alpha),
\]
since $\alpha - H\alpha = \Delta G\alpha = d\delta G\alpha + \delta dG\alpha$. So the Green operator, an inverse of $\Delta$ on the orthogonal complement of the harmonic forms, produces the primitives: if $\alpha$ is
exact, then $\alpha = d(\delta G\alpha)$, and $\delta G\alpha$ is the primitive of $\alpha$ orthogonal to the closed forms, the primitive of smallest norm.

\begin{example}\label{hodge-theorem:example-circle-decomposition}
On the circle $M = \RR/2\pi\ZZ$ with the coordinate $\theta$, write functions as Fourier series $f = \sum_kc_ke^{ik\theta}$ (complex coefficients, $c_{-k} = \bar c_k$ for real $f$). Here $\Delta f = -f^{\prime\prime}$,
and on $1$-forms $\delta(g\,d\theta) = -g^{\prime}$ (by Definition~\ref{riemannian-geometry:definition-laplacian}, $\delta = -{*}d{*}$ on $1$-forms when $n = 1$, and $*(g\,d\theta) = g$). The harmonic functions are the
constants and the harmonic $1$-forms the constant multiples of $d\theta$, so the decompositions of Theorem~\ref{hodge-theorem:theorem-decomposition} are
\[
f = c_0 + \delta\Bigl(-\sum_{k \neq 0}\frac{c_k}{ik}e^{ik\theta}\,d\theta\Bigr), \qquad g\,d\theta = c_0\,d\theta + d\Bigl(\sum_{k \neq 0}\frac{c_k}{ik}e^{ik\theta}\Bigr) \quad\text{for } g = \sum_kc_ke^{ik\theta},
\]
and the Green operator on functions is $G(\sum_kc_ke^{ik\theta}) = \sum_{k \neq 0}c_ke^{ik\theta}/k^2$. The harmonic part of $g\,d\theta$ is $c_0\,d\theta = \bigl(\frac{1}{2\pi}\int_0^{2\pi}g\bigr)d\theta$, and the theorem
says that $g\,d\theta$ is exact exactly when its average vanishes, which is the familiar condition for $\int_0^\theta g$ to be periodic.
\end{example}

\begin{example}\label{hodge-theorem:example-helmholtz}
On the flat torus $T^3 = \RR^3/\ZZ^3$ identify $1$-forms with vector fields by the metric, $X = (X_1, X_2, X_3) \leftrightarrow X_1\,dx_1 + X_2\,dx_2 + X_3\,dx_3$, and $2$-forms with vector fields by
$Y \leftrightarrow *Y^\flat = Y_1\,dx_2 \wedge dx_3 + Y_2\,dx_3 \wedge dx_1 + Y_3\,dx_1 \wedge dx_2$. Then $df \leftrightarrow \nabla f$, the codifferential of a $2$-form is $\delta(*Y^\flat) = *d{*}{*}Y^\flat = *dY^\flat$
(here $n = 3$, $k = 2$, so the sign in $\delta = (-1)^{n(k+1)+1}*d*$ is $+1$, and $** = 1$), and $*dY^\flat \leftrightarrow \operatorname{curl}Y$ (Example~\ref{differential-forms:example-forms-r3}). The harmonic
$1$-forms are the constant ones (Example~\ref{riemannian-geometry:example-harmonic-forms}). So Theorem~\ref{hodge-theorem:theorem-decomposition} in degree $1$ is the \emph{Helmholtz decomposition}:
every smooth vector field $X$ on $T^3$ can be written as
\[
X = c + \nabla f + \operatorname{curl}Y
\]
with $c$ a constant vector, the three summands are uniquely determined by $X$ (although $f$ and $Y$ are not), and they are orthogonal in $L^2$; moreover $\operatorname{curl}X = 0$ if and only if $X = c + \nabla f$ (the closed $1$-forms), and $\operatorname{div}X = 0$ if
and only if $X = c + \operatorname{curl}Y$ (the coclosed ones, as $\delta X^\flat = -\operatorname{div}X$). In Fourier coefficients, $X = \sum_k\hat X(k)e^{2\pi ik\cdot x}$, the three parts of the term with
$k \neq 0$ are the components of $\hat X(k)$ along $k$ and orthogonal to $k$, and $c = \hat X(0)$.
\end{example}

\begin{theorem}[Hodge]\label{hodge-theorem:theorem-hodge}
The map $\mathcal{H}^k \to H^k_{\mathrm{dR}}(M)$, $\alpha \mapsto [\alpha]$, is an isomorphism: every de Rham class contains a unique harmonic representative, namely the one of smallest $L^2$ norm. In
particular $H^k_{\mathrm{dR}}(M)$ is finite-dimensional, and $\dim\mathcal{H}^k = b_k(M)$.
\end{theorem}

\begin{proof}
Harmonic forms are closed, so the map is defined. It is injective: if $\alpha \in \mathcal{H}^k$ is exact, $\alpha = d\beta$, then $\|\alpha\|^2 = (\alpha, d\beta) = (\delta\alpha, \beta) = 0$. It is surjective: a closed
$\alpha$ decomposes as $\alpha = \eta + d\beta + \delta\gamma$ with $\eta$ harmonic, and $\delta\gamma = 0$ as in the proof of Theorem~\ref{hodge-theorem:theorem-decomposition}, so $\alpha = \eta + d\beta$ is
cohomologous to $\eta$. If $\alpha = \eta + d\beta$ is any representative, then $\|\alpha\|^2 = \|\eta\|^2 + \|d\beta\|^2 \geq \|\eta\|^2$ with equality only for $\alpha = \eta$, which is the statement about the
norm. Finite-dimensionality and the identification with the Betti numbers follow from Theorem~\ref{de-rham:theorem-de-rham}.
\end{proof}

The minimisation property has a variational reading: among all closed forms in a class, the harmonic one minimises the energy $\|\alpha\|^2$, and the condition $\delta\alpha = 0$ is the
Euler--Lagrange equation, since $\frac{d}{dt}\|\alpha + t\,d\beta\|^2\big|_{t=0} = 2(\alpha, d\beta) = 2(\delta\alpha, \beta)$ must vanish for all $\beta$.

\begin{corollary}\label{hodge-theorem:corollary-poincare-duality}
The Hodge star induces an isomorphism $* \colon \mathcal{H}^k \to \mathcal{H}^{n-k}$, hence $b_k(M) = b_{n-k}(M)$, and the pairing
\[
H^k_{\mathrm{dR}}(M) \times H^{n-k}_{\mathrm{dR}}(M) \to \RR, \qquad ([\alpha], [\beta]) \mapsto \int_M\alpha \wedge \beta,
\]
is nondegenerate; for $[\beta] = [*\alpha]$ with $\alpha$ harmonic and nonzero it takes the value $\|\alpha\|^2 > 0$.
\end{corollary}

\begin{proof}
$*$ commutes with $\Delta$ (Proposition~\ref{riemannian-geometry:proposition-adjoint}(4)) and is an isomorphism of $\Omega^k$ onto $\Omega^{n-k}$, hence restricts to an isomorphism on harmonic forms; combine
with Theorem~\ref{hodge-theorem:theorem-hodge}. For a nonzero harmonic $\alpha$, $\int_M\alpha \wedge *\alpha = \|\alpha\|^2 > 0$, so the pairing is nondegenerate on the harmonic representatives, which represent
all classes.
\end{proof}

\begin{remark}\label{hodge-theorem:remark-dependence-on-metric}
The space $\mathcal{H}^k$ depends on the metric, but its dimension does not, and the isomorphism of Theorem~\ref{hodge-theorem:theorem-hodge} identifies all of them with $H^k_{\mathrm{dR}}(M)$. This is the
mechanism that produces the Hodge structures of Chapter~\ref{kahler}: on a compact K\"ahler manifold the harmonic forms decompose by type, and the resulting decomposition of cohomology is
independent of the metric, although the harmonic representatives are not (Exercise~\ref{hodge-theorem:exercise-harmonic-not-metric-independent}).
\end{remark}

On a surface part of this dependence disappears: harmonic $1$-forms only see the conformal class of the metric. This is the first sign of the relation between harmonic forms on a surface and its
structure of Riemann surface.

\begin{proposition}\label{hodge-theorem:proposition-conformal-invariance}
Let $M$ be an oriented surface and $g^{\prime} = e^{2u}g$ a conformal metric. Then the Hodge stars of $g$ and $g^{\prime}$ agree on $1$-forms, so a $1$-form is harmonic for $g$ if and only if it is harmonic
for $g^{\prime}$ (when $M$ is compact).
\end{proposition}

\begin{proof}
If $(\omega^1, \omega^2)$ is an oriented orthonormal coframe for $g$, then $(e^u\omega^1, e^u\omega^2)$ is one for $g^{\prime}$. The Hodge star of $g$ maps $\omega^1 \mapsto \omega^2$ and $\omega^2 \mapsto -\omega^1$
(Theorem~\ref{multilinear-algebra:theorem-hodge-star}), and that of $g^{\prime}$ maps $e^u\omega^1 \mapsto e^u\omega^2$ and $e^u\omega^2 \mapsto -e^u\omega^1$; by linearity over functions the two stars agree on
$\omega^1$ and $\omega^2$, hence on all $1$-forms. On a compact surface a $1$-form $\alpha$ is harmonic if and only if $d\alpha = 0$ and $d{*}\alpha = 0$
(Proposition~\ref{riemannian-geometry:proposition-adjoint}(3), as $\delta = -{*}d{*}$ on $1$-forms in dimension $2$), and these conditions only involve $*$ on $1$-forms.
\end{proof}

For a compact Riemann surface, with any metric compatible with its complex structure, $*$ on $1$-forms is determined by the complex structure ($*\,dx = dy$ and $*\,dy = -dx$ in a holomorphic coordinate
$z = x + iy$), and a real $1$-form $\alpha = a\,dx + b\,dy$ is harmonic exactly when $\alpha + i{*}\alpha = (a - ib)\,dz$ is a holomorphic $1$-form (the Cauchy--Riemann equations for $a - ib$ are
$a_y = b_x$ and $a_x + b_y = 0$, that is, $d\alpha = 0$ and $d{*}\alpha = 0$); this identifies $\mathcal{H}^1$ with the real parts of holomorphic $1$-forms, and is the
case of curves of the Hodge decomposition $H^1 = H^{1,0} \oplus H^{0,1}$ (Chapter~\ref{kahler}).

\begin{proposition}\label{hodge-theorem:proposition-lie-group-harmonic}
Let $G$ be a compact connected Lie group with a bi-invariant Riemannian metric (Example~\ref{riemannian-geometry:example-metrics}) and an orientation. Then the harmonic forms on $G$ are exactly the bi-invariant forms.
\end{proposition}

\begin{proof}
Left and right translations are isometries, and they preserve the orientation, being connected to the identity through diffeomorphisms; so they commute with the Hodge star, and $*$ maps bi-invariant forms to bi-invariant forms. A bi-invariant
form $\omega$ is closed (Theorem~\ref{de-rham:theorem-lie-group-cohomology}), and so is $*\omega$; hence $\delta\omega = \pm{*}d{*}\omega = 0$ and $\omega$ is harmonic (Proposition~\ref{riemannian-geometry:proposition-adjoint}(3)). The
bi-invariant $k$-forms and the harmonic $k$-forms both map isomorphically onto $H^k_{\mathrm{dR}}(G)$ (Theorems~\ref{de-rham:theorem-lie-group-cohomology} and~\ref{hodge-theorem:theorem-hodge}), so the inclusion of the first
space into the second is an equality.
\end{proof}

\begin{counterexample}\label{hodge-theorem:counterexample-compactness}
Compactness is essential in the Hodge theorem. On $\RR$ with the Euclidean metric, the harmonic functions are the affine functions $ax + b$, a two-dimensional space, while $H^0_{\mathrm{dR}}(\RR) = \RR$; and the harmonic $1$-forms $(ax + b)\,dx$
form a two-dimensional space although $H^1_{\mathrm{dR}}(\RR) = 0$. Without compactness harmonic forms need not be closed or coclosed: the function $x$ is harmonic with $dx \neq 0$, and $x\,dx$ is harmonic with $\delta(x\,dx) = -1 \neq 0$, since $\delta = -{*}d{*}$ on $1$-forms on $\RR$. The
integration by parts behind Proposition~\ref{riemannian-geometry:proposition-adjoint}(3) fails for forms without compact support.
\end{counterexample}

\section{Applications}\label{hodge-theorem:section-applications}

\begin{corollary}\label{hodge-theorem:corollary-euler-characteristic}
For a compact oriented manifold, $\sum_k(-1)^k\dim\mathcal{H}^k = \chi(M)$, which is the index of the elliptic operator
$d + \delta \colon \Omega^{\mathrm{even}}(M) \to \Omega^{\mathrm{odd}}(M)$.
\end{corollary}

\begin{proof}
The first equality is Theorem~\ref{hodge-theorem:theorem-hodge} together with the definition of $\chi$ as the alternating sum of Betti numbers
(Corollary~\ref{poincare-duality:corollary-euler-odd}). For the second, $(d + \delta)^* = d + \delta$ exchanges even and odd forms, its kernel on even forms is $\bigoplus_{k \text{ even}}\mathcal{H}^k$ and on odd
forms $\bigoplus_{k \text{ odd}}\mathcal{H}^k$, because $(d + \delta)\alpha = 0$ implies $\Delta\alpha = (d+\delta)^2\alpha = 0$; so the index is $\sum_k(-1)^k\dim\mathcal{H}^k$.
\end{proof}

\begin{corollary}[Poincar\'e inequality]\label{hodge-theorem:corollary-poincare-inequality}
Let $M$ be compact and connected, and let $\lambda_1 > 0$ be the smallest nonzero eigenvalue of the Laplacian $\Delta = \delta d$ on functions. Then for every smooth function $f$ with $\int_Mf\,\mathrm{vol} = 0$,
\[
\int_Mf^2\,\mathrm{vol} \leq \frac{1}{\lambda_1}\int_M|df|^2\,\mathrm{vol},
\]
with equality exactly for the eigenfunctions of eigenvalue $\lambda_1$. (For intervals, slabs and cubes in $\RR^n$ see Subsection~\ref{distributions:subsection-poincare}.)
\end{corollary}

\begin{proof}
A function $u$ with $\Delta u = 0$ satisfies $\|du\|^2 = (\delta du, u) = 0$, so it is locally constant, hence constant as $M$ is connected: $\mathcal{H}^0 = \RR$. By Theorem~\ref{elliptic-operators:theorem-spectral} there is an
orthonormal basis $u_0, u_1, \ldots$ of $L^2(M)$ of smooth eigenfunctions with eigenvalues $0 = \lambda^{(0)} < \lambda_1 = \lambda^{(1)} \leq \lambda^{(2)} \leq \cdots$, where $u_0$ is constant. Write
$f = \sum_kc_ku_k$; then $c_0 = (f, u_0) = 0$ by hypothesis. Since $\Delta f = \sum_k\lambda^{(k)}c_ku_k$ in $L^2$,
\[
\|df\|^2 = (\Delta f, f) = \sum_{k \geq 1}\lambda^{(k)}c_k^2 \geq \lambda_1\sum_{k \geq 1}c_k^2 = \lambda_1\|f\|^2,
\]
with equality if and only if $c_k = 0$ whenever $\lambda^{(k)} > \lambda_1$, that is, $f$ lies in the $\lambda_1$-eigenspace.
\end{proof}

\begin{example}\label{hodge-theorem:example-wirtinger}
On the circle $M = \RR/2\pi\ZZ$ with its standard metric, $\Delta f = -f''$. The real solutions of $f'' = -\lambda f$ with $\lambda > 0$ are $a\cos(\sqrt\lambda\,\theta) + b\sin(\sqrt\lambda\,\theta)$, which are
$2\pi$-periodic only for $\sqrt\lambda \in \ZZ$; so the eigenvalues are the $k^2$, with eigenfunctions $\cos k\theta$, $\sin k\theta$ (and $1$ for $k = 0$). Hence $\lambda_1 = 1$, and
Corollary~\ref{hodge-theorem:corollary-poincare-inequality} gives back Wirtinger's inequality (Proposition~\ref{distributions:proposition-wirtinger} with $L = 2\pi$), with its extremals
$a\cos\theta + b\sin\theta$; on a general compact manifold the corollary is its replacement, with the constant $1/\lambda_1$. Theorem~\ref{elliptic-operators:theorem-spectral} also shows that the functions $1$, $\cos k\theta$, $\sin k\theta$ span a dense subspace of $L^2$, which is the completeness of Fourier series.
\end{example}

The last application relates harmonic $1$-forms to curvature. It rests on a formula expressing the Hodge Laplacian through the Levi-Civita connection. Below $\nabla$ also denotes the connections
induced on $1$-forms and $2$-forms (Proposition~\ref{connections:proposition-induced-connections}), so that $(\nabla_X\alpha)(Y) = X(\alpha(Y)) - \alpha(\nabla_XY)$ and
$(\nabla_X\beta)(Y, Z) = X(\beta(Y, Z)) - \beta(\nabla_XY, Z) - \beta(Y, \nabla_XZ)$, and $(e_i)$ is a local orthonormal frame. By Lemma~\ref{characteristic-classes:lemma-normal-frames}, applied to the
metric connection $\nabla$ on $TM$, near every point $p$ there are orthonormal frames with $\nabla e_i = 0$ at $p$; then also $[e_i, e_j] = \nabla_{e_i}e_j - \nabla_{e_j}e_i = 0$ at $p$, and the dual coframe
$(e^i)$ satisfies $\nabla e^i = 0$ at $p$, its connection form being $-\theta^T$.

\begin{lemma}\label{hodge-theorem:lemma-codifferential-connection}
Let $\alpha$ be a $1$-form and $\beta$ a $2$-form on $M$.
\begin{enumerate}
\item $d\alpha(X, Y) = (\nabla_X\alpha)(Y) - (\nabla_Y\alpha)(X)$.
\item For a vector field $V$, $d(\iota_V\mathrm{vol}) = (\operatorname{div}V)\,\mathrm{vol}$ with $\operatorname{div}V = \sum_i\langle\nabla_{e_i}V, e_i\rangle$; hence $\int_M\operatorname{div}V\,\mathrm{vol} = 0$.
\item $\delta\alpha = -\sum_i(\nabla_{e_i}\alpha)(e_i)$ and $(\delta\beta)(X) = -\sum_i(\nabla_{e_i}\beta)(e_i, X)$.
\end{enumerate}
\end{lemma}

\begin{proof}
(1) By the invariant formula (Theorem~\ref{differential-forms:theorem-exterior-derivative}) and the Leibniz rule for $\nabla$,
$d\alpha(X, Y) = X(\alpha(Y)) - Y(\alpha(X)) - \alpha([X, Y]) = (\nabla_X\alpha)(Y) + \alpha(\nabla_XY) - (\nabla_Y\alpha)(X) - \alpha(\nabla_YX) - \alpha([X, Y])$, and the terms without derivatives of $\alpha$ cancel
because $\nabla_XY - \nabla_YX = [X, Y]$.

(2) The trace $\sum_i\langle\nabla_{e_i}V, e_i\rangle$ does not depend on the orthonormal frame, so we may compute at a point $p$ in a frame with $\nabla e_i(p) = 0$, oriented, with dual coframe $(e^i)$. Then
$\mathrm{vol} = e^1 \wedge \cdots \wedge e^n$ and, with $V = \sum_iV^ie_i$, $\iota_V\mathrm{vol} = \sum_i(-1)^{i-1}V^i\,e^1 \wedge \cdots \widehat{e^i} \cdots \wedge e^n$. By (1), $de^j(X, Y) = (\nabla_Xe^j)(Y) - (\nabla_Ye^j)(X)$ vanishes at
$p$, so at $p$ only the derivatives of the $V^i$ contribute: $d(\iota_V\mathrm{vol}) = \sum_i(-1)^{i-1}dV^i \wedge e^1 \wedge \cdots \widehat{e^i} \cdots \wedge e^n = \sum_ie_i(V^i)\,\mathrm{vol}$. And at $p$,
$e_i(V^i) = e_i\langle V, e_i\rangle = \langle\nabla_{e_i}V, e_i\rangle$. The integral vanishes by Stokes' theorem.

(3) Let $\delta^{\prime}$ be defined by the right-hand sides; again they do not depend on the frame. We show that $\delta^{\prime}$ is a formal adjoint of $d$; then $(\gamma, \delta\beta - \delta^{\prime}\beta) = 0$ for
all $1$-forms $\gamma$, so $\delta\beta = \delta^{\prime}\beta$, and similarly in degree $1$. Let $\gamma$ be a $1$-form, and $V$ the vector field with $\langle V, X\rangle = \sum_j\gamma(e_j)\beta(X, e_j)$. At a point $p$, in a frame with
$\nabla e_i(p) = 0$, the derivative of a function $\gamma(e_j)$ along $e_i$ is $(\nabla_{e_i}\gamma)(e_j)$, and similarly for $\beta$, so
\[
\operatorname{div}V = \sum_{i,j}e_i\bigl(\gamma(e_j)\beta(e_i, e_j)\bigr) = \sum_{i,j}(\nabla_{e_i}\gamma)(e_j)\beta(e_i, e_j) + \sum_{i,j}\gamma(e_j)(\nabla_{e_i}\beta)(e_i, e_j) = \langle d\gamma, \beta\rangle - \langle\gamma, \delta^{\prime}\beta\rangle,
\]
where we used $\langle d\gamma, \beta\rangle = \frac12\sum_{i,j}d\gamma(e_i, e_j)\beta(e_i, e_j) = \sum_{i,j}(\nabla_{e_i}\gamma)(e_j)\beta(e_i, e_j)$, by (1) and the antisymmetry of $\beta$. Integrating with (2),
$(d\gamma, \beta) = (\gamma, \delta^{\prime}\beta)$. For a function $f$ and a $1$-form $\alpha$ the same computation with $V = f\alpha^\sharp$, where $\langle\alpha^\sharp, X\rangle = \alpha(X)$, gives
$\operatorname{div}V = \sum_ie_i(f)\alpha(e_i) + f\sum_i(\nabla_{e_i}\alpha)(e_i) = \langle df, \alpha\rangle - f\,\delta^{\prime}\alpha$.
\end{proof}

\begin{proposition}[Bochner--Weitzenb\"ock formula]\label{hodge-theorem:proposition-weitzenbock}
For a $1$-form $\alpha$ on $M$,
\[
\Delta\alpha = \nabla^*\nabla\alpha + \mathrm{Ric}(\alpha^\sharp, -), \qquad \nabla^*\nabla\alpha = -\sum_i\bigl(\nabla_{e_i}\nabla_{e_i}\alpha - \nabla_{\nabla_{e_i}e_i}\alpha\bigr),
\]
and $(\Delta\alpha, \alpha) = \|\nabla\alpha\|^2 + \int_M\mathrm{Ric}(\alpha^\sharp, \alpha^\sharp)\,\mathrm{vol}$, where $\|\nabla\alpha\|^2 = \int_M\sum_i|\nabla_{e_i}\alpha|^2\,\mathrm{vol}$.
\end{proposition}

\begin{proof}
All terms are independent of the frame, so we compute at a point $p$ in a frame with $\nabla e_i = 0$ and $[e_i, e_j] = 0$ at $p$; then at $p$ the derivative along $e_i$ of a function $\phi(e_j)$ is
$(\nabla_{e_i}\phi)(e_j)$ for any form $\phi$. By Lemma~\ref{hodge-theorem:lemma-codifferential-connection}(3), $\delta\alpha = -\sum_i(\nabla_{e_i}\alpha)(e_i)$ near $p$, so at $p$
\[
(d\delta\alpha)(e_j) = e_j(\delta\alpha) = -\sum_i(\nabla_{e_j}\nabla_{e_i}\alpha)(e_i).
\]
By the same lemma and (1), $(\delta d\alpha)(e_j) = -\sum_ie_i\bigl(d\alpha(e_i, e_j)\bigr) = -\sum_i(\nabla_{e_i}\nabla_{e_i}\alpha)(e_j) + \sum_i(\nabla_{e_i}\nabla_{e_j}\alpha)(e_i)$ at $p$. Adding,
\[
(\Delta\alpha)(e_j) = -\sum_i(\nabla_{e_i}\nabla_{e_i}\alpha)(e_j) + \sum_i\bigl((\nabla_{e_i}\nabla_{e_j} - \nabla_{e_j}\nabla_{e_i})\alpha\bigr)(e_i) .
\]
The first sum is $(\nabla^*\nabla\alpha)(e_j)$ at $p$. Since $[e_i, e_j] = 0$ at $p$, the operator in the second sum is the curvature $R^*(e_i, e_j)$ of the connection on $1$-forms, and
$(R^*(X, Y)\alpha)(Z) = -\alpha(R(X, Y)Z)$: applying $XY - YX - [X, Y]$, which is zero on functions, to $\alpha(Z)$ and expanding with the Leibniz rule, the terms $(\nabla_X\alpha)(\nabla_YZ) + (\nabla_Y\alpha)(\nabla_XZ)$
cancel and what remains is $(R^*(X, Y)\alpha)(Z) + \alpha(R(X, Y)Z) = 0$. With $A = \alpha^\sharp$ the second sum is
\[
-\sum_i\langle A, R(e_i, e_j)e_i\rangle = -\sum_iR(e_i, e_j, e_i, A) = \sum_iR(e_i, e_j, A, e_i) = \mathrm{Ric}(e_j, A),
\]
by the antisymmetry of $R$ in its last two arguments (Proposition~\ref{riemannian-geometry:proposition-curvature-symmetries}) and Definition~\ref{riemannian-geometry:definition-curvature-tensor}. For the integral formula, at $p$
\[
\langle\nabla^*\nabla\alpha, \alpha\rangle = -\sum_i\langle\nabla_{e_i}\nabla_{e_i}\alpha, \alpha\rangle = \sum_i|\nabla_{e_i}\alpha|^2 - \sum_ie_i\langle\nabla_{e_i}\alpha, \alpha\rangle = \sum_i|\nabla_{e_i}\alpha|^2 - \operatorname{div}W
\]
with $W = \sum_i\langle\nabla_{e_i}\alpha, \alpha\rangle e_i$, the vector field with $\langle W, X\rangle = \langle\nabla_X\alpha, \alpha\rangle = \frac12X|\alpha|^2$. Integrating with
Lemma~\ref{hodge-theorem:lemma-codifferential-connection}(2) gives $(\nabla^*\nabla\alpha, \alpha) = \|\nabla\alpha\|^2$, and $\langle\mathrm{Ric}(A, -), \alpha\rangle = \mathrm{Ric}(A, A)$.
\end{proof}

\begin{theorem}[Bochner]\label{hodge-theorem:theorem-bochner}
Let $(M, g)$ be a compact oriented Riemannian manifold with $\mathrm{Ric} \geq 0$. Then every harmonic $1$-form is parallel, and $b_1(M) \leq n$. If moreover $\mathrm{Ric} > 0$ at some point, then
$b_1(M) = 0$.
\end{theorem}

\begin{proof}
By Proposition~\ref{hodge-theorem:proposition-weitzenbock}, for $\alpha$ harmonic,
\[
0 = (\Delta\alpha, \alpha) = \|\nabla\alpha\|^2 + \int_M\mathrm{Ric}(\alpha^\sharp, \alpha^\sharp)\,\mathrm{vol},
\]
where $\alpha^\sharp$ is the vector field dual to $\alpha$. Both terms are non-negative, hence both vanish: $\nabla\alpha = 0$, so $\alpha$ is parallel and determined by its value at one point, giving
$\dim\mathcal{H}^1 \leq n$; and $\mathrm{Ric}(\alpha^\sharp, \alpha^\sharp) = 0$ everywhere. If $\mathrm{Ric} > 0$ at $p$ and $\alpha \neq 0$, then $\alpha$ is nowhere zero (being parallel and nonzero somewhere), so
$\mathrm{Ric}(\alpha^\sharp, \alpha^\sharp)(p) > 0$, a contradiction; hence $\mathcal{H}^1 = 0$ and $b_1(M) = 0$ by Theorem~\ref{hodge-theorem:theorem-hodge}.
\end{proof}

\begin{example}\label{hodge-theorem:example-applications}
\begin{enumerate}
\item On the round sphere $S^n$ with $n \geq 2$, $\mathrm{Ric} > 0$, so $b_1 = 0$, which agrees with Example~\ref{de-rham:example-spheres}.
\item On a flat torus the harmonic forms are those with constant coefficients (Example~\ref{riemannian-geometry:example-harmonic-forms}), so Theorem~\ref{hodge-theorem:theorem-hodge} recovers
$b_k(T^n) = \binom{n}{k}$.
\item A compact oriented manifold with a metric of positive Ricci curvature has $b_1 = 0$; by Bochner's argument applied to the other Laplacians, positivity assumptions on curvature produce vanishing
theorems in all degrees. In the complex setting the same mechanism gives the Kodaira vanishing theorem (Chapter~\ref{kodaira}).
\end{enumerate}
\end{example}

Both bounds of Theorem~\ref{hodge-theorem:theorem-bochner} are sharp. The flat torus $T^n$ has $\mathrm{Ric} = 0$ and $b_1 = n$, its harmonic $1$-forms $dx_i$ being parallel. The product $S^1 \times S^{n-1}$
($n \geq 3$) of a circle and a round sphere has $\mathrm{Ric} \geq 0$, with $\mathrm{Ric}(\partial_\theta, \partial_\theta) = 0$ along the circle factor, and $b_1 = 1$: the harmonic form $d\theta$ is parallel, and
$\mathrm{Ric}(\partial_\theta, \partial_\theta) = 0$ is exactly what the proof requires.

\begin{remark}\label{hodge-theorem:remark-elliptic-complexes}
The proof of Theorem~\ref{hodge-theorem:theorem-decomposition} used only that $(\Omega^\bullet, d)$ is a complex of sections of vector bundles whose Laplacian is elliptic: an \emph{elliptic complex}. The same
argument applies verbatim to any elliptic complex, in particular to the Dolbeault complex $(\Omega^{p,\bullet}, \bar\partial)$ on a compact complex manifold, giving the isomorphism between Dolbeault
cohomology and $\bar\partial$-harmonic forms (Chapter~\ref{complex-manifolds}), which is how the Hodge decomposition of Chapter~\ref{kahler} is proved. There is also a proof of the Hodge theorem by the
heat equation, which constructs the harmonic projection as $\lim_{t \to \infty}e^{-t\Delta}$ and yields the local index formula \cite{lawson-michelsohn}; for functions this limit is computed in
Exercise~\ref{elliptic-operators:exercise-heat-equation}.
\end{remark}

\section{Elliptic complexes}\label{hodge-theorem:section-elliptic-complexes}

The de Rham complex is not the only complex of differential operators to which the argument of Section~\ref{hodge-theorem:section-decomposition} applies. What the proof used is that the Laplacian
$d\delta + \delta d$ is elliptic, although $d$ itself is not; and the ellipticity of the Laplacian is a pointwise statement about the symbols. We isolate this property; it is the setting of the Dolbeault
complex of a complex manifold (Chapter~\ref{complex-manifolds}), of de Rham complexes with twisted coefficients, and of the index theorem.

\begin{definition}\label{hodge-theorem:definition-elliptic-complex}
Let $M$ be a compact manifold and $E^0, \ldots, E^N$ vector bundles over $M$ with Hermitian (or, for real bundles, Euclidean) metrics. A \emph{complex of differential operators} is a sequence
\[
0 \to \Gamma(M, E^0) \xrightarrow{D_0} \Gamma(M, E^1) \xrightarrow{D_1} \cdots \xrightarrow{D_{N-1}} \Gamma(M, E^N) \to 0
\]
of differential operators $D_k$ of the same order $m \geq 1$ with $D_{k+1} \circ D_k = 0$. It is an \emph{elliptic complex} if for every $x \in M$ and every $\xi \in T^*_xM \setminus \{0\}$ the sequence of principal symbols
\[
0 \to E^0_x \xrightarrow{\sigma_m(D_0)(x, \xi)} E^1_x \xrightarrow{\sigma_m(D_1)(x, \xi)} \cdots \to E^N_x \to 0
\]
is exact. Its \emph{cohomology} is $H^k(E^\bullet, D) = \ker D_k/\operatorname{im}D_{k-1}$.
\end{definition}

\begin{notation}\label{hodge-theorem:notation-elliptic-complex}
We write $D$ for each of the operators $D_k$ and $D^*$ for their formal adjoints (Definition~\ref{elliptic-operators:definition-formal-adjoint}), so that $D^* \colon \Gamma(E^{k+1}) \to \Gamma(E^k)$; the
\emph{Laplacians} of the complex are $\Delta_k = D_{k-1}D_{k-1}^* + D_k^*D_k \colon \Gamma(E^k) \to \Gamma(E^k)$, often simply written $\Delta = DD^* + D^*D$, and
$\mathcal{H}^k(E^\bullet) = \ker\Delta_k$ is the space of \emph{harmonic sections}. Sections of $E^k$ are \emph{closed} if $Du = 0$ and \emph{coclosed} if $D^*u = 0$.
\end{notation}

The symbols of a complex form a complex: by the multiplicativity of principal symbols (Proposition~\ref{elliptic-operators:proposition-symbol}), $\sigma_m(D_{k+1})\sigma_m(D_k) = \sigma_{2m}(D_{k+1}D_k) = 0$.
Ellipticity requires this complex of finite-dimensional vector spaces to be exact. For $N = 1$ it says that the single operator $D_0 \colon \Gamma(E^0) \to \Gamma(E^1)$ is elliptic in the sense of
Definition~\ref{elliptic-operators:definition-differential-operator}, so elliptic complexes generalise elliptic operators.

\begin{lemma}\label{hodge-theorem:lemma-laplacian-elliptic-complex}
The Laplacians of an elliptic complex are elliptic and formally self-adjoint operators of order $2m$, and $\int_M\langle\Delta u, u\rangle = \|Du\|^2 + \|D^*u\|^2$ for smooth $u$.
\end{lemma}

\begin{proof}
\emph{Step 1: the symbol.} By the multiplicativity of symbols and $\sigma_m(D^*) = \sigma_m(D)^*$ (Definition~\ref{elliptic-operators:definition-formal-adjoint}), writing $a_k = \sigma_m(D_k)(x, \xi)$,
\[
\sigma_{2m}(\Delta_k)(x, \xi) = a_{k-1}a_{k-1}^* + a_k^*a_k .
\]
\emph{Step 2: the finite-dimensional lemma.} The maps $a_k$ form a complex of finite-dimensional inner product spaces, and the right-hand side is its Laplacian in degree $k$. The proof of
Lemma~\ref{hodge-theorem:lemma-finite-dimensional} works verbatim for complex inner product spaces, and shows that the kernel of this Laplacian is isomorphic to the cohomology of the complex
$(a_\bullet)$ in degree $k$. For $\xi \neq 0$ this cohomology vanishes by ellipticity, so $\sigma_{2m}(\Delta_k)(x, \xi)$ is injective, hence invertible. So $\Delta_k$ is elliptic.

\emph{Step 3: adjointness.} $\Delta_k^* = (D_{k-1}D_{k-1}^*)^* + (D_k^*D_k)^* = \Delta_k$, and $\int\langle\Delta u, u\rangle = \int\langle D^*u, D^*u\rangle + \int\langle Du, Du\rangle$ by the definition of the formal adjoint.
\end{proof}

\begin{theorem}[Hodge theorem for elliptic complexes]\label{hodge-theorem:theorem-elliptic-complex}
Let $(\Gamma(E^\bullet), D)$ be an elliptic complex over a compact manifold. For each $k$:
\begin{enumerate}
\item $\mathcal{H}^k(E^\bullet)$ is finite-dimensional, consists of smooth sections, and equals $\ker D \cap \ker D^*$;
\item there is an orthogonal decomposition $\Gamma(E^k) = \mathcal{H}^k(E^\bullet) \oplus D\Gamma(E^{k-1}) \oplus D^*\Gamma(E^{k+1})$, with $\ker D_k = \mathcal{H}^k(E^\bullet) \oplus D\Gamma(E^{k-1})$;
\item the map $\mathcal{H}^k(E^\bullet) \to H^k(E^\bullet, D)$, $u \mapsto [u]$, is an isomorphism; in particular the cohomology of an elliptic complex is finite-dimensional.
\end{enumerate}
\end{theorem}

\begin{proof}
\emph{Step 1: harmonic sections.} By Lemma~\ref{hodge-theorem:lemma-laplacian-elliptic-complex} and Theorem~\ref{elliptic-operators:theorem-fredholm}(1), $\ker\Delta_k$ is finite-dimensional and consists of smooth
sections; and $\Delta u = 0$ gives $\|Du\|^2 + \|D^*u\|^2 = 0$, so $\ker\Delta = \ker D \cap \ker D^*$, the converse being clear.

\emph{Step 2: decomposition.} By Corollary~\ref{elliptic-operators:corollary-p-star-p} applied to $\Delta_k = \Delta_k^*$, $\Gamma(E^k) = \Delta\Gamma(E^k) \oplus \ker\Delta$, and
$\Delta\Gamma(E^k) \subseteq D\Gamma(E^{k-1}) + D^*\Gamma(E^{k+1})$. The three spaces are pairwise orthogonal: for $h$ harmonic, $\int\langle h, Dv\rangle = \int\langle D^*h, v\rangle = 0$ and
$\int\langle h, D^*w\rangle = \int\langle Dh, w\rangle = 0$; and $\int\langle Dv, D^*w\rangle = \int\langle DDv, w\rangle = 0$. So the sum is direct and orthogonal. If $Du = 0$ and $u = h + Dv + D^*w$, then
$0 = \int\langle Du, w\rangle = \int\langle DD^*w, w\rangle = \|D^*w\|^2$, so $u \in \mathcal{H}^k \oplus D\Gamma(E^{k-1})$; the reverse inclusion is clear since $DD = 0$.

\emph{Step 3: cohomology.} Harmonic sections are closed. By Step 2 every closed section is $h + Dv$ with $h$ harmonic, which gives surjectivity, and if $h = Dv$ is harmonic then
$\|h\|^2 = \int\langle h, Dv\rangle = \int\langle D^*h, v\rangle = 0$, which gives injectivity.
\end{proof}

\begin{example}\label{hodge-theorem:example-elliptic-complexes}
\begin{enumerate}
\item \emph{The de Rham complex.} Here $E^k = \Lambda^kT^*M$, $D = d$ and $\sigma_1(d)(x, \xi) = i\xi \wedge -$ (Example~\ref{elliptic-operators:example-symbol-computations}(5)). The sequence
$0 \to \Lambda^0 \xrightarrow{\xi\wedge} \Lambda^1 \xrightarrow{\xi\wedge} \cdots$ is exact for $\xi \neq 0$: choose a basis $e^1 = \xi, e^2, \ldots, e^n$ of $T^*_xM$ and write $\alpha = e^1 \wedge \beta + \gamma$ with $\beta$,
$\gamma$ not involving $e^1$; then $\xi \wedge \alpha = e^1 \wedge \gamma$, which vanishes only if $\gamma = 0$, that is, if $\alpha = \xi \wedge \beta$. So
Theorem~\ref{hodge-theorem:theorem-elliptic-complex} contains Theorems~\ref{hodge-theorem:theorem-decomposition} and~\ref{hodge-theorem:theorem-hodge}.
\item \emph{Twisted de Rham complexes.} Let $E \to M$ be a vector bundle with a flat connection $\nabla$ (Chapter~\ref{connections}) and any Hermitian metric. The exterior covariant derivative
$d_\nabla \colon \Omega^k(M, E) \to \Omega^{k+1}(M, E)$ satisfies $d_\nabla \circ d_\nabla = 0$ because the curvature vanishes, and its symbol is $i\xi \wedge - \otimes \id_E$, since in a local frame
$d_\nabla = d + A\wedge$ differs from $d$ by an operator of order $0$. The symbol sequence is the one of (1) tensored with $E_x$, hence exact. So the cohomology of a compact manifold with coefficients in a
flat bundle (a local system) is finite-dimensional and is represented by harmonic forms; the metric on $E$ need not be parallel, but then $d_\nabla^*$ is not simply $\pm{*}d_\nabla{*}$.
\item \emph{The Dolbeault complex} $(\Omega^{p,\bullet}(X), \bar\partial)$ of a compact complex manifold is elliptic; this is the case used in Chapter~\ref{complex-manifolds}.
\item \emph{A single operator}: for $N = 1$, Theorem~\ref{hodge-theorem:theorem-elliptic-complex} says that $\ker D$ and $\Gamma(E^1)/D\Gamma(E^0) \cong \ker D^*$ are finite-dimensional, which is
part of Theorem~\ref{elliptic-operators:theorem-fredholm} (with Corollary~\ref{elliptic-operators:corollary-p-star-p}).
\end{enumerate}
\end{example}

\begin{counterexample}\label{hodge-theorem:counterexample-not-elliptic-complex}
Exactness of the symbol sequence cannot be dropped. On the torus $T^2 = \RR^2/\ZZ^2$ consider the complex $0 \to C^\infty(T^2) \xrightarrow{\partial_x} C^\infty(T^2) \to 0$. Its symbol $i\xi_1$ vanishes at
$\xi = dy$, so it is not elliptic, and indeed its cohomology $H^0 = \ker\partial_x$ contains every smooth function of $y$ alone, an infinite-dimensional space. Likewise $H^1 = C^\infty(T^2)/\partial_xC^\infty(T^2)$
is infinite-dimensional: $f \mapsto \int_0^1f(x, y)\,dx$ vanishes on $\partial_xC^\infty$ and maps onto $C^\infty(\RR/\ZZ)$.
\end{counterexample}

\begin{corollary}\label{hodge-theorem:corollary-euler-characteristic-complex}
For an elliptic complex, the operator $D + D^* \colon \bigoplus_{k \text{ even}}\Gamma(E^k) \to \bigoplus_{k \text{ odd}}\Gamma(E^k)$ is elliptic, and its index equals the \emph{Euler characteristic of the complex}
$\chi(E^\bullet, D) = \sum_k(-1)^k\dim H^k(E^\bullet, D)$.
\end{corollary}

\begin{proof}
On $\bigoplus_k\Gamma(E^k)$ the operator $P = D + D^*$ is formally self-adjoint and $P^2 = DD^* + D^*D = \Delta$, since $DD = 0$ and $D^*D^* = 0$. So $\sigma_m(P)^2 = \sigma_{2m}(\Delta)$ is invertible for $\xi \neq 0$, hence
so is $\sigma_m(P)$, and its restriction from even to odd degrees is an isomorphism (it is injective on each, and $\sigma_m(P)$ exchanges the two). The kernel of $P$ on even sections is
$\bigoplus_{k \text{ even}}\mathcal{H}^k$: $Pu = 0$ gives $\Delta u = P^2u = 0$, and conversely harmonic sections satisfy $Du = D^*u = 0$. The cokernel is the kernel of the adjoint, $P$ on odd sections, which is
$\bigoplus_{k \text{ odd}}\mathcal{H}^k$ for the same reason. Conclude with Theorem~\ref{hodge-theorem:theorem-elliptic-complex}(3).
\end{proof}

For the de Rham complex this is Corollary~\ref{hodge-theorem:corollary-euler-characteristic}. The index of an elliptic complex is computed from characteristic classes by the Atiyah--Singer index theorem
(Chapter~\ref{spin-geometry}); the Euler characteristic and the signature of a manifold, and the holomorphic Euler characteristic of a complex manifold, are the three classical instances.

\section{The spectrum of the Hodge Laplacian}\label{hodge-theorem:section-spectrum}

The Hodge theorem concerns the eigenvalue $0$ of the Laplacian. The positive eigenvalues carry no topology individually, but they come in a pattern that the topology constrains, and their size is constrained by the
curvature.

\begin{notation}\label{hodge-theorem:notation-eigenspaces}
For $\lambda \in \RR$ let $E^k_\lambda = \{\alpha \in \Omega^k(M) : \Delta\alpha = \lambda\alpha\}$. By Theorem~\ref{elliptic-operators:theorem-spectral}, $E^k_\lambda$ is finite-dimensional, nonzero only for countably many
$\lambda \geq 0$, and $E^k_0 = \mathcal{H}^k$. We write $\lambda_1^{(k)}$ for the smallest positive eigenvalue of $\Delta$ on $k$-forms.
\end{notation}

\begin{proposition}\label{hodge-theorem:proposition-eigenspace-complex}
Let $\lambda > 0$. Then $d(E^k_\lambda) \subseteq E^{k+1}_\lambda$ and $\delta(E^k_\lambda) \subseteq E^{k-1}_\lambda$, the complex
\[
0 \to E^0_\lambda \xrightarrow{d} E^1_\lambda \xrightarrow{d} \cdots \xrightarrow{d} E^n_\lambda \to 0
\]
is exact, $E^k_\lambda = dE^{k-1}_\lambda \oplus \delta E^{k+1}_\lambda$ orthogonally, and $d$ restricts to an isomorphism $\delta E^{k+1}_\lambda \to dE^k_\lambda$. Consequently $\sum_k(-1)^k\dim E^k_\lambda = 0$.
\end{proposition}

\begin{proof}
\emph{Step 1: invariance.} $d$ and $\delta$ commute with $\Delta$ (Proposition~\ref{riemannian-geometry:proposition-adjoint}), so they map $\lambda$-eigenforms to $\lambda$-eigenforms.

\emph{Step 2: decomposition.} For $\alpha \in E^k_\lambda$, $\alpha = \lambda^{-1}\Delta\alpha = d(\lambda^{-1}\delta\alpha) + \delta(\lambda^{-1}d\alpha)$, with $\lambda^{-1}\delta\alpha \in E^{k-1}_\lambda$ and $\lambda^{-1}d\alpha \in E^{k+1}_\lambda$ by
Step 1. The two summands are orthogonal, as in Theorem~\ref{hodge-theorem:theorem-decomposition}.

\emph{Step 3: exactness.} If $\alpha \in E^k_\lambda$ and $d\alpha = 0$, Step 2 gives $\alpha = d(\lambda^{-1}\delta\alpha)$, so $\alpha$ is exact within the complex. For the isomorphism, $d$ maps $\delta E^{k+1}_\lambda$ onto
$d(dE^{k-1}_\lambda \oplus \delta E^{k+1}_\lambda) = dE^k_\lambda$, and it is injective there: if $d\delta\beta = 0$ then $\|\delta\beta\|^2 = (\beta, d\delta\beta) = 0$.

\emph{Step 4: alternating sum.} An exact sequence of finite-dimensional vector spaces has vanishing alternating sum of dimensions (by rank--nullity, $\dim E^k_\lambda = \dim\ker d_k + \dim\operatorname{im}d_k$
and $\ker d_k = \operatorname{im}d_{k-1}$, so the sum telescopes).
\end{proof}

\begin{corollary}[McKean--Singer identity, finite form]\label{hodge-theorem:corollary-mckean-singer}
For every $\Lambda \geq 0$ and every $t \in \RR$,
\[
\sum_{k=0}^n(-1)^k\sum_{\lambda \leq \Lambda}e^{-t\lambda}\dim E^k_\lambda = \chi(M).
\]
\end{corollary}

\begin{proof}
The terms with $\lambda > 0$ cancel by Proposition~\ref{hodge-theorem:proposition-eigenspace-complex}, and the terms with $\lambda = 0$ give $\sum_k(-1)^k\dim\mathcal{H}^k = \chi(M)$ by
Corollary~\ref{hodge-theorem:corollary-euler-characteristic}.
\end{proof}

\begin{remark}\label{hodge-theorem:remark-heat-kernel}
Letting $\Lambda \to \infty$, the left-hand side becomes $\sum_k(-1)^k\operatorname{tr}e^{-t\Delta_k}$; the series converge because the eigenvalues grow like $k^{2/n}$ (Weyl's law; on tori this is
Proposition~\ref{elliptic-operators:proposition-weyl-torus}). So the supertrace of the heat operator is independent of $t$. As $t \to \infty$ it tends to $\sum_k(-1)^k\dim\mathcal{H}^k$, which is the Hodge
theorem; as $t \to 0$ the heat kernel has an asymptotic expansion whose coefficients are integrals of local curvature expressions, and comparing the two limits proves the Gauss--Bonnet--Chern theorem and,
for twisted Dirac operators, the index theorem \cite{lawson-michelsohn}.
\end{remark}

\begin{example}\label{hodge-theorem:example-circle-eigenspaces}
On $M = \RR/2\pi\ZZ$, $E^0_{j^2}$ is spanned by $\cos j\theta$ and $\sin j\theta$ and $E^1_{j^2}$ by $\cos j\theta\,d\theta$ and $\sin j\theta\,d\theta$ for $j \geq 1$ (Example~\ref{hodge-theorem:example-wirtinger}). The map
$d \colon E^0_{j^2} \to E^1_{j^2}$ sends $\cos j\theta \mapsto -j\sin j\theta\,d\theta$ and $\sin j\theta \mapsto j\cos j\theta\,d\theta$, an isomorphism, as Proposition~\ref{hodge-theorem:proposition-eigenspace-complex} predicts;
and $\chi(S^1) = 1 - 1 = 0$.
\end{example}

The eigenvalue expansion also gives the heat equation proof of the Hodge theorem, announced in Remark~\ref{hodge-theorem:remark-elliptic-complexes}: the heat flow deforms a closed form inside its cohomology
class to the harmonic representative.

\begin{proposition}[Heat flow]\label{hodge-theorem:proposition-heat-flow}
Let $\alpha \in \Omega^k(M)$, let $(u_j)$ be an orthonormal basis of eigenforms of $\Delta$ on $k$-forms with eigenvalues $\lambda_j$ (Theorem~\ref{elliptic-operators:theorem-spectral}), and
$\alpha(t) = \sum_je^{-t\lambda_j}(\alpha, u_j)u_j$ for $t \geq 0$. Then:
\begin{enumerate}
\item $\alpha(t)$ is smooth in $(t, x)$, $\partial_t\alpha + \Delta\alpha = 0$ and $\alpha(0) = \alpha$;
\item $\|\alpha(t) - H\alpha\| \leq e^{-t\lambda_1^{(k)}}\|\alpha - H\alpha\|$, and $\alpha(t) \to H\alpha$ in every $C^r$ norm as $t \to \infty$;
\item $d(\alpha(t)) = (d\alpha)(t)$; if $\alpha$ is closed, every $\alpha(t)$ is closed and cohomologous to $\alpha$, and the energy $\|\alpha(t)\|^2$ decreases to $\|H\alpha\|^2$.
\end{enumerate}
\end{proposition}

\begin{proof}
(1) The coefficients $c_j = (\alpha, u_j)$ decay faster than any power of $\lambda_j$ (proof of Theorem~\ref{elliptic-operators:theorem-spectral}), and $|e^{-t\lambda_j}| \leq 1$, so the series and its term-by-term
derivatives in $t$ converge in every Sobolev norm, uniformly in $t \geq 0$; hence $\alpha(t)$ is smooth and the equation holds term by term.

(2) $H\alpha = \sum_{\lambda_j = 0}c_ju_j$, so $\|\alpha(t) - H\alpha\|^2 = \sum_{\lambda_j > 0}e^{-2t\lambda_j}c_j^2 \leq e^{-2t\lambda_1^{(k)}}\|\alpha - H\alpha\|^2$. The same estimate holds for
$\Delta^N(\alpha(t) - H\alpha) = (\Delta^N\alpha)(t) - H\Delta^N\alpha$, since $(\Delta^N\alpha, u_j) = \lambda_j^Nc_j$. The operator $\Delta^N$ is elliptic of order $2N$ (its symbol is $|\xi|^{2N}$), so
Theorem~\ref{elliptic-operators:theorem-elliptic-estimate} bounds the $H^{2N}$ norm of $\alpha(t) - H\alpha$ by a constant times $e^{-t\lambda_1^{(k)}}(\|\Delta^N\alpha\| + \|\alpha\|)$. By the Sobolev embedding, used as in the proof of Theorem~\ref{elliptic-operators:theorem-spectral}, the $H^{2N}$ norm bounds the $C^r$ norm for
$N$ large.

(3) Let $v$ be an eigenform of $\Delta$ on $(k+1)$-forms, with eigenvalue $\mu$. Then $\delta v \in E^k_\mu$, so $(u_j, \delta v) = 0$ unless $\lambda_j = \mu$ (eigenspaces of a symmetric operator for different
eigenvalues are orthogonal), and
\[\begin{aligned}
(d(\alpha(t)), v) &= (\alpha(t), \delta v) = \sum_je^{-t\lambda_j}c_j(u_j, \delta v)\\
&= e^{-t\mu}(\alpha, \delta v) = e^{-t\mu}(d\alpha, v) = ((d\alpha)(t), v).
\end{aligned}\]
As the eigenforms $v$ are complete in $L^2$, $d(\alpha(t)) = (d\alpha)(t)$. If $d\alpha = 0$ then $d\alpha(t) = 0$, and
$\alpha(t) - \alpha = -\int_0^t\Delta\alpha(s)\,ds = d\bigl(-\int_0^t\delta\alpha(s)\,ds\bigr)$ is exact. Finally $\|\alpha(t)\|^2 = \sum_je^{-2t\lambda_j}c_j^2$ is non-increasing with limit $\sum_{\lambda_j = 0}c_j^2 = \|H\alpha\|^2$.
\end{proof}

The size of the first eigenvalue on functions is controlled by the Ricci curvature. The next theorem, due to Lichnerowicz, is the second application of the Bochner--Weitzenb\"ock formula.

\begin{theorem}[Lichnerowicz]\label{hodge-theorem:theorem-lichnerowicz}
Let $(M, g)$ be a compact oriented Riemannian manifold of dimension $n \geq 2$ with $\mathrm{Ric} \geq (n - 1)\kappa\,g$ for a constant $\kappa > 0$. Then the smallest positive eigenvalue of the Laplacian on functions
satisfies $\lambda_1^{(0)} \geq n\kappa$.
\end{theorem}

\begin{proof}
Let $\Delta f = \lambda f$ with $\lambda > 0$ and $f \neq 0$, and $\alpha = df$.

\emph{Step 1: the two sides of the Bochner formula.} Since $d$ commutes with $\Delta$, $\Delta\alpha = d\Delta f = \lambda\alpha$, and $\|\alpha\|^2 = (\delta df, f) = \lambda\|f\|^2$. So
Proposition~\ref{hodge-theorem:proposition-weitzenbock} gives
\[\begin{aligned}
\lambda^2\|f\|^2 = \lambda\|\alpha\|^2 = (\Delta\alpha, \alpha) &= \|\nabla\alpha\|^2 + \int_M\mathrm{Ric}(\alpha^\sharp, \alpha^\sharp)\,\mathrm{vol}\\
&\geq \|\nabla\alpha\|^2 + (n - 1)\kappa\lambda\|f\|^2 .
\end{aligned}\]
\emph{Step 2: the Hessian term.} In a local orthonormal frame let $h_{ij} = (\nabla_{e_i}\alpha)(e_j)$, so that $|\nabla\alpha|^2 = \sum_{i,j}h_{ij}^2$ and, by Lemma~\ref{hodge-theorem:lemma-codifferential-connection}(3),
$\Delta f = \delta\alpha = -\sum_ih_{ii}$. By the Cauchy--Schwarz inequality, $(\sum_ih_{ii})^2 \leq n\sum_ih_{ii}^2 \leq n\sum_{i,j}h_{ij}^2$. Hence $|\nabla\alpha|^2 \geq \frac1n(\Delta f)^2$ pointwise, and
$\|\nabla\alpha\|^2 \geq \frac1n\|\Delta f\|^2 = \frac{\lambda^2}{n}\|f\|^2$.

\emph{Step 3: conclusion.} Combining, $\lambda^2 \geq \frac{\lambda^2}{n} + (n - 1)\kappa\lambda$; dividing by $\lambda > 0$, $\lambda\frac{n - 1}{n} \geq (n - 1)\kappa$, that is, $\lambda \geq n\kappa$.
\end{proof}

\begin{example}\label{hodge-theorem:example-lichnerowicz-sphere}
The bound is attained on the round unit sphere $S^n$, where $\mathrm{Ric} = (n - 1)g$ (so $\kappa = 1$): the restriction of a linear function $z$ to $S^n$ satisfies $\Delta z = nz$
(Exercise~\ref{hodge-theorem:exercise-linear-functions-sphere}), so $\lambda_1^{(0)} \leq n$, and Theorem~\ref{hodge-theorem:theorem-lichnerowicz} gives $\lambda_1^{(0)} = n$. For $n = 2$ this is the value $2$ of
Exercise~\ref{elliptic-operators:exercise-first-eigenvalue-sphere}, now obtained without knowing the whole spectrum.
\end{example}

\begin{counterexample}\label{hodge-theorem:counterexample-lichnerowicz}
A positive lower bound on the Ricci curvature is needed. The flat torus $\RR^n/L\ZZ^n$ has $\mathrm{Ric} = 0$ and first eigenvalue $4\pi^2/L^2$, attained by $\cos(2\pi x_1/L)$, which tends to $0$ as $L \to \infty$; so no
bound depending only on $n$ holds when $\kappa = 0$. And a lower bound on $\mathrm{Ric}$ cannot be replaced by a lower bound on the scalar curvature: the product $S^2(1) \times S^1(L)$ of a unit sphere and a circle of
length $L$ has scalar curvature $2 > 0$ but first eigenvalue at most $4\pi^2/L^2$, again by a function of the circle coordinate.
\end{counterexample}

\section{Harmonic forms in the middle degree}\label{hodge-theorem:section-middle-degree}

Proposition~\ref{hodge-theorem:proposition-conformal-invariance} has an analogue in every even dimension: the Hodge star on forms of half the dimension depends only on the conformal class. In dimensions divisible by four
this leads to a description of the signature by harmonic forms.

\begin{proposition}\label{hodge-theorem:proposition-conformal-star}
Let $g^{\prime} = e^{2u}g$ be conformal metrics on an oriented $n$-manifold, with Hodge stars $*$ and $*^{\prime}$. Then $*^{\prime} = e^{(n - 2k)u}\,*$ on $k$-forms. In particular, for $n = 2m$, $*^{\prime} = *$ on $m$-forms, and
on a compact manifold of dimension $2m$ the space $\mathcal{H}^m$ depends only on the conformal class of the metric.
\end{proposition}

\begin{proof}
Let $(\omega^1, \ldots, \omega^n)$ be an oriented orthonormal coframe for $g$; then $(e^u\omega^1, \ldots, e^u\omega^n)$ is one for $g^{\prime}$. For a multi-index $I$ of length $k$ with complement $I^c$, the
Hodge star of $g$ maps $\omega^I \mapsto \epsilon_I\omega^{I^c}$ with a sign $\epsilon_I$ depending only on the permutation $(I, I^c)$ (Theorem~\ref{multilinear-algebra:theorem-hodge-star}), and that of $g^{\prime}$ maps
$e^{ku}\omega^I \mapsto \epsilon_Ie^{(n-k)u}\omega^{I^c}$. By linearity over functions, $*^{\prime}\omega^I = e^{(n-2k)u}*\omega^I$. On a compact manifold an $m$-form is harmonic if and only if $d\alpha = 0$ and
$d{*}\alpha = 0$ (Proposition~\ref{riemannian-geometry:proposition-adjoint}(3) and $\delta = \pm{*}d{*}$), conditions which, for $n = 2m$, involve the metric only through $*$ on $m$-forms.
\end{proof}

\begin{definition}\label{hodge-theorem:definition-self-dual}
Let $M$ be an oriented Riemannian manifold of dimension $n = 4l$. On $2l$-forms, $** = (-1)^{2l \cdot 2l} = 1$ (Theorem~\ref{multilinear-algebra:theorem-hodge-star}), so $\Lambda^{2l}T^*M = \Lambda^+ \oplus \Lambda^-$ with
$\Lambda^{\pm} = \{\alpha : *\alpha = \pm\alpha\}$. Forms with $*\alpha = \alpha$ are \emph{self-dual}, those with $*\alpha = -\alpha$ \emph{anti-self-dual}. When $M$ is compact, $\mathcal{H}^\pm$ denotes the space of
harmonic self-dual, respectively anti-self-dual, $2l$-forms, and $b^\pm = \dim\mathcal{H}^\pm$.
\end{definition}

Since $*$ is an isometry with $*^2 = 1$ on $\Lambda^{2l}$, it is self-adjoint there, and $\Lambda^+ \perp \Lambda^-$ pointwise; every $2l$-form splits as $\alpha = \frac12(\alpha + *\alpha) + \frac12(\alpha - *\alpha)$.
As $*$ commutes with $\Delta$, it preserves $\mathcal{H}^{2l}$, and $\mathcal{H}^{2l} = \mathcal{H}^+ \oplus \mathcal{H}^-$.

\begin{theorem}\label{hodge-theorem:theorem-signature-harmonic}
Let $M$ be a compact oriented Riemannian manifold of dimension $4l$. Then the signature of $M$ (Definition~\ref{poincare-duality:definition-intersection-form}) is $\sigma(M) = b^+ - b^-$, and $b^+ + b^- = b_{2l}(M)$.
\end{theorem}

\begin{proof}
\emph{Step 1: the form on harmonic representatives.} Under the de Rham isomorphism the cup product corresponds to the wedge product (Theorem~\ref{de-rham:theorem-de-rham}), so
$\sigma(M)$ is the signature of $Q(\alpha, \beta) = \int_M\alpha \wedge \beta$ on $H^{2l}_{\mathrm{dR}}(M)$, which by Theorem~\ref{hodge-theorem:theorem-hodge} we may compute on $\mathcal{H}^{2l} = \mathcal{H}^+ \oplus \mathcal{H}^-$.

\emph{Step 2: definiteness.} For $\alpha \in \mathcal{H}^\pm$, $Q(\alpha, \alpha) = \int\alpha \wedge \alpha = \pm\int\alpha \wedge *\alpha = \pm\|\alpha\|^2$. So $Q$ is positive definite on $\mathcal{H}^+$ and negative definite on $\mathcal{H}^-$.

\emph{Step 3: orthogonality.} For $\alpha \in \mathcal{H}^+$ and $\beta \in \mathcal{H}^-$, $Q(\alpha, \beta) = \int\alpha \wedge \beta = \int\beta \wedge *\alpha = (\beta, \alpha) = 0$, using that forms of even degree commute and that
$\mathcal{H}^+ \perp \mathcal{H}^-$. So in a basis adapted to $\mathcal{H}^+ \oplus \mathcal{H}^-$ the matrix of $Q$ is block diagonal, positive definite of size $b^+$ and negative definite of size $b^-$, and
$\sigma(M) = b^+ - b^-$ by Sylvester's law of inertia. The last statement is Theorem~\ref{hodge-theorem:theorem-hodge}.
\end{proof}

\begin{example}\label{hodge-theorem:example-self-dual}
\begin{enumerate}
\item On $\RR^4$ with the standard metric and orientation, $*(dx_1 \wedge dx_2) = dx_3 \wedge dx_4$, $*(dx_1 \wedge dx_3) = -dx_2 \wedge dx_4$, $*(dx_1 \wedge dx_4) = dx_2 \wedge dx_3$, so $\Lambda^+$ has the basis
$dx_1 \wedge dx_2 + dx_3 \wedge dx_4$, $dx_1 \wedge dx_3 - dx_2 \wedge dx_4$, $dx_1 \wedge dx_4 + dx_2 \wedge dx_3$, and $\Lambda^-$ the basis with the opposite signs: both have rank $3$.
\item On the flat torus $T^4 = \RR^4/\ZZ^4$ the harmonic $2$-forms are those with constant coefficients (Example~\ref{riemannian-geometry:example-harmonic-forms}), so by (1) $b^+ = b^- = 3$ and $\sigma(T^4) = 0$.
\item On $S^2 \times S^2$ with the product of two unit round metrics, let $\omega_1$, $\omega_2$ be the pullbacks of the two area forms. They are closed and parallel, and $*\omega_1 = \omega_2$, since
$\omega_1 \wedge \omega_2$ is the volume form and $\omega_1$, $\omega_2$ are built from complementary halves of an oriented orthonormal coframe. Hence $\omega_1$ and $\omega_2$ are harmonic, they span
$\mathcal{H}^2$ (as $b_2 = 2$ by Theorem~\ref{de-rham:theorem-kunneth}), $\mathcal{H}^+ = \RR(\omega_1 + \omega_2)$ and $\mathcal{H}^- = \RR(\omega_1 - \omega_2)$, and $\sigma = 0$, as in
Example~\ref{poincare-duality:example-intersection-forms}.
\end{enumerate}
\end{example}

\begin{remark}\label{hodge-theorem:remark-signature-operator}
On all forms of a $4l$-manifold there is an involution $\tau$, equal to $*$ up to a sign depending on the degree, which anticommutes with $d + \delta$. The operator $d + \delta$ from the $(+1)$-eigenspace to the
$(-1)$-eigenspace of $\tau$ is the \emph{signature operator}; its index is $\sigma(M)$, by the argument of Theorem~\ref{hodge-theorem:theorem-signature-harmonic} (the contributions of the degrees other than $2l$ cancel in
pairs). The Atiyah--Singer theorem computes this index as the $L$-genus, which is Hirzebruch's signature theorem (Remark~\ref{characteristic-classes:remark-signature}). In dimension $4$ the self-dual and anti-self-dual
$2$-forms are the starting point of Donaldson theory: a connection whose curvature is anti-self-dual (an \emph{instanton}) minimises the Yang--Mills functional.
\end{remark}

\section{Products}\label{hodge-theorem:section-products}

\begin{proposition}\label{hodge-theorem:proposition-products}
Let $M$ and $N$ be compact oriented Riemannian manifolds and $M \times N$ the product with the product metric and orientation, with projections $p$ and $q$. If $\alpha \in \mathcal{H}^a(M)$ and $\beta \in \mathcal{H}^b(N)$,
then $p^*\alpha \wedge q^*\beta$ is harmonic, and the map
\[
\bigoplus_{a+b=k}\mathcal{H}^a(M) \otimes \mathcal{H}^b(N) \to \mathcal{H}^k(M \times N), \qquad \alpha \otimes \beta \mapsto p^*\alpha \wedge q^*\beta,
\]
is an isomorphism. In particular $\mathcal{H}^k(M \times N)$ depends only on $\mathcal{H}^\bullet(M)$ and $\mathcal{H}^\bullet(N)$.
\end{proposition}

\begin{proof}
\emph{Step 1: the star of a product form.} If $(\omega^1, \ldots, \omega^m)$ and $(\theta^1, \ldots, \theta^{n})$ are oriented orthonormal coframes on $M$ and $N$, their pullbacks form an oriented orthonormal coframe on
$M \times N$. For multi-indices $I$, $J$ of lengths $a$, $b$, with $\omega^I \wedge \omega^{I^c} = \epsilon_I\omega^{1\cdots m}$ and similarly for $J$, one checks on these monomials that
\[
*(p^*\omega^I \wedge q^*\theta^J) = (-1)^{b(m - a)}\,p^*(*\omega^I) \wedge q^*(*\theta^J),
\]
since $(\omega^I \wedge \theta^J) \wedge (\omega^{I^c} \wedge \theta^{J^c}) = (-1)^{b(m-a)}(\omega^I \wedge \omega^{I^c}) \wedge (\theta^J \wedge \theta^{J^c})$. By linearity over functions of $p$ and of $q$ the formula holds for all
$\alpha \in \Omega^a(M)$, $\beta \in \Omega^b(N)$.

\emph{Step 2: harmonicity.} If $\alpha$ and $\beta$ are harmonic, then $\alpha$, $*\alpha$, $\beta$ and $*\beta$ are closed, so $p^*\alpha \wedge q^*\beta$ and, by Step 1, its Hodge star are closed. By
Proposition~\ref{riemannian-geometry:proposition-adjoint}(3) the product is harmonic.

\emph{Step 3: bijectivity.} The map is the composition of the isomorphisms of Theorem~\ref{hodge-theorem:theorem-hodge} for $M$ and $N$, the K\"unneth isomorphism of Theorem~\ref{de-rham:theorem-kunneth} (compact
manifolds have finite good covers \cite[Chapter I, Section 5]{bott-tu}), and the inverse of the isomorphism of Theorem~\ref{hodge-theorem:theorem-hodge} for $M \times N$, since the harmonic form
$p^*\alpha \wedge q^*\beta$ represents the class $[p^*\alpha \wedge q^*\beta]$.
\end{proof}

\begin{example}\label{hodge-theorem:example-products}
On $S^1 \times S^{n-1}$ with the product of round metrics, $\mathcal{H}^\bullet(S^1)$ is spanned by $1$ and $d\theta$, and $\mathcal{H}^\bullet(S^{n-1})$ by $1$ and the volume form $\mathrm{vol}_{S^{n-1}}$ (Exercise~\ref{hodge-theorem:exercise-harmonic-top}
and Example~\ref{de-rham:example-spheres}). So the harmonic forms on the product are spanned by $1$, $d\theta$, $\mathrm{vol}_{S^{n-1}}$ and $d\theta \wedge \mathrm{vol}_{S^{n-1}}$; for $n \geq 3$ this agrees with the
discussion of the sharpness of Theorem~\ref{hodge-theorem:theorem-bochner}, where $d\theta$ appeared as the parallel harmonic $1$-form.
\end{example}

\section{Exercises}\label{hodge-theorem:section-exercises}

\begin{exercise}\label{hodge-theorem:exercise-harmonic-top}
Show that on a compact connected oriented Riemannian $n$-manifold, $\mathcal{H}^0 = \RR$ and $\mathcal{H}^n = \RR\,\mathrm{vol}_g$, and deduce $b_0 = b_n = 1$.
\end{exercise}

\begin{solution}
A harmonic function has $df = 0$, hence is constant; apply $*$ for the top degree (Corollary~\ref{hodge-theorem:corollary-poincare-duality}).
\end{solution}

\begin{exercise}\label{hodge-theorem:exercise-green-properties}
Show that the Green operator is self-adjoint, that it is compact as an operator on $L^2$, and that $\alpha \mapsto H\alpha$ is the orthogonal projection onto $\mathcal{H}^k$ with $H = I - \Delta G$.
\end{exercise}

\begin{solution}
$H = I - \Delta G$ is the identity $\alpha = H\alpha + \Delta G\alpha$ of Theorem~\ref{hodge-theorem:theorem-decomposition}. For self-adjointness write $\alpha = H\alpha + \Delta a$ and $\beta = H\beta + \Delta b$
with $a = G\alpha$, $b = G\beta$ orthogonal to $\mathcal{H}^k$. Then $(G\alpha, \beta) = (a, H\beta + \Delta b) = (a, \Delta b)$, since $a \perp \mathcal{H}^k$, and likewise $(\alpha, G\beta) = (\Delta a, b)$; these are equal because
$\Delta$ is formally self-adjoint. For compactness, use an orthonormal basis of eigenforms $u_k$ of $\Delta$ (Theorem~\ref{elliptic-operators:theorem-spectral}): $G$ kills the $u_k$ with eigenvalue $0$ and maps
$u_k$ to $\lambda_k^{-1}u_k$ otherwise. As $\lambda_k \to \infty$, $G$ is the norm limit of the finite rank operators obtained by truncating this diagonal operator, so it is compact
(Proposition~\ref{functional-analysis:proposition-compact-operators}), and it extends from smooth forms to $L^2$.
\end{solution}

\begin{exercise}\label{hodge-theorem:exercise-harmonic-not-metric-independent}
Give an example of a compact manifold and two metrics for which the harmonic representatives of a fixed cohomology class differ. Why can the two metrics not be conformal to each other when the
manifold is a surface and the class has degree $1$?
\end{exercise}

\begin{solution}
On the torus $T^2 = \RR^2/\ZZ^2$ take the flat metric $dx^2 + dy^2$, for which the harmonic representative of $[dx]$ is $dx$, and the metric $g^{\prime} = e^{2\phi(x)}dx^2 + dy^2$ with $\phi$ smooth,
$1$-periodic and non-constant. An oriented orthonormal coframe for $g^{\prime}$ is $(e^\phi dx, dy)$, so its Hodge star maps $e^\phi dx \mapsto dy$ and $dy \mapsto -e^\phi dx$. The form
$\alpha = c\,e^{\phi(x)}dx$ with $c = (\int_0^1e^\phi\,dx)^{-1}$ is closed (its coefficient depends only on $x$) and $*^{\prime}\alpha = c\,dy$ is closed, so $\alpha$ is $g^{\prime}$-harmonic; its integrals over the
two circles $y = 0$ and $x = 0$ are $1$ and $0$, as for $dx$, so it lies in the class $[dx]$ (on the torus a class in $H^1$ is determined by these two integrals, Exercise~\ref{de-rham:exercise-torus}).
It differs from $dx$ because $\phi$ is not constant; equivalently, $dx$ is not $g^{\prime}$-harmonic, since $*^{\prime}dx = e^{-\phi}dy$ has $d(e^{-\phi}dy) = -\phi^{\prime}e^{-\phi}dx \wedge dy \neq 0$. For conformal metrics on a
surface this cannot happen, by Proposition~\ref{hodge-theorem:proposition-conformal-invariance}; and $g^{\prime}$ is indeed not conformal to the flat metric, as $dx$ and $dy$ have different $g^{\prime}$-lengths where
$\phi \neq 0$, while they remain orthogonal.
\end{solution}

\begin{exercise}\label{hodge-theorem:exercise-hodge-riemann-surface}
Let $M$ be a compact orientable surface of genus $g$. Using Theorem~\ref{hodge-theorem:theorem-hodge} and
Corollary~\ref{hodge-theorem:corollary-poincare-duality}, show that $\dim\mathcal{H}^1 = 2g$ and that the intersection pairing on $\mathcal{H}^1$ is a nondegenerate alternating form.
\end{exercise}

\begin{solution}
$\dim\mathcal{H}^1 = b_1(M) = 2g$ by Theorem~\ref{hodge-theorem:theorem-hodge} and the computation of the Betti numbers of $\Sigma_g$ (Example~\ref{singular-homology:example-cellular-computations}). The pairing
$(\alpha, \beta) \mapsto \int_M\alpha \wedge \beta$ is bilinear and alternating, since $\alpha \wedge \beta = -\beta \wedge \alpha$ for $1$-forms, and it is nondegenerate on $\mathcal{H}^1$: for $\alpha \neq 0$
harmonic, $*\alpha$ is harmonic (Corollary~\ref{hodge-theorem:corollary-poincare-duality}) and $\int_M\alpha \wedge *\alpha = \|\alpha\|^2 > 0$. (A nondegenerate alternating form forces the dimension to be even,
in accordance with $2g$.)
\end{solution}

\begin{exercise}\label{hodge-theorem:exercise-finite-dimensional}
Let $V^0 = \RR$, $V^1 = \RR^2$, $V^2 = \RR$ with the standard inner products, $d \colon V^0 \to V^1$, $t \mapsto (t, t)$, and $d = 0 \colon V^1 \to V^2$. Compute the cohomology, the
Laplacians and their kernels in each degree, and verify Lemma~\ref{hodge-theorem:lemma-finite-dimensional}. Which element of the cohomology class of $(1, 0)$ has the smallest norm?
\end{exercise}

\begin{solution}
$H^0 = \ker(t \mapsto (t, t)) = 0$, $H^1 = \RR^2/\RR(1, 1) \cong \RR$ and $H^2 = V^2 = \RR$. The adjoint of $d \colon V^0 \to V^1$ is $\delta(a, b) = a + b$, and the
adjoint of the zero map is zero. So $\Delta = \delta d = 2$ on $V^0$, with kernel $0$; on $V^1$, $\Delta(a, b) = d\delta(a, b) = (a + b, a + b)$, with kernel the line $a + b = 0$; and $\Delta = 0$ on $V^2$, with kernel
$V^2$. The dimensions $0$, $1$, $1$ of the kernels are those of the cohomology. The class of $(1, 0)$ is $(1, 0) + \RR(1, 1)$, and $(1, 0) = (\frac12, -\frac12) + d(\frac12)$ with $(\frac12, -\frac12)$ harmonic: it
is the element of the class of smallest norm, $\frac{1}{\sqrt2} < 1$, the orthogonal projection of $(1, 0)$ to the line $a + b = 0$.
\end{solution}

\begin{exercise}\label{hodge-theorem:exercise-weitzenbock-sphere}
Verify the Bochner--Weitzenb\"ock formula for $\alpha = dz|_{S^2}$ on the round unit sphere $S^2 \subseteq \RR^3$. Use that $\Delta z = 2z$ (Exercise~\ref{elliptic-operators:exercise-spherical-harmonics}), so
$\Delta\alpha = d\Delta z = 2\alpha$, and compute $\nabla^*\nabla\alpha$ from the Hessian $\nabla dz = -z\,g$ of the restriction of a linear function to the sphere.
\end{exercise}

\begin{solution}
For a function $f$, $\nabla df$ is the Hessian $(X, Y) \mapsto (\nabla_Xdf)(Y) = X(Y(f)) - (\nabla_XY)(f)$. For the restriction of $z$ and tangent fields $X$, $Y$ on $S^2$, $Y(z) = \langle Y, e_3\rangle$
because $z$ is linear, so $X(Y(z)) = \langle\overline\nabla_XY, e_3\rangle$, where $\overline\nabla$ is the derivative in $\RR^3$. By the Gauss formula (Theorem~\ref{riemannian-geometry:theorem-gauss-formula}),
$\overline\nabla_XY = \nabla_XY + \mathrm{II}(X, Y)$, and on the unit sphere the normal part is $\mathrm{II}(X, Y) = \langle\overline\nabla_XY, x\rangle x = -\langle Y, \overline\nabla_Xx\rangle x = -\langle X, Y\rangle x$ at the point $x$.
Hence $(\nabla_Xdz)(Y) = \langle\overline\nabla_XY - \nabla_XY, e_3\rangle = -\langle X, Y\rangle\langle x, e_3\rangle = -\langle X, Y\rangle z$. Hence $\nabla_{e_i}\alpha = -z\,e^i$ in an orthonormal frame, and
$\nabla^*\nabla\alpha = -\sum_i\nabla_{e_i}(-z\,e^i)$, which at a point with $\nabla e_i = 0$ is $\sum_ie_i(z)\,e^i = dz = \alpha$. Since $\mathrm{Ric} = g$ on the unit sphere, the formula reads
$\Delta\alpha = \alpha + \alpha = 2\alpha$, as it should.
\end{solution}

\begin{exercise}\label{hodge-theorem:exercise-linear-functions-sphere}
Let $z$ be the restriction to the round unit sphere $S^n \subseteq \RR^{n+1}$ of a linear function. Show that $\Delta z = nz$, and deduce that $n$ is an eigenvalue of the Laplacian on functions with multiplicity at least
$n + 1$.
\end{exercise}

\begin{solution}
The computation in the solution of Exercise~\ref{hodge-theorem:exercise-weitzenbock-sphere} works in every dimension: with $z(x) = \langle x, e\rangle$, the Gauss formula gives $(\nabla_Xdz)(Y) = -\langle X, Y\rangle z$ for
tangent fields $X$, $Y$. By Lemma~\ref{hodge-theorem:lemma-codifferential-connection}(3), $\Delta z = \delta dz = -\sum_{i=1}^n(\nabla_{e_i}dz)(e_i) = nz$. The restrictions of the $n + 1$ coordinate functions are linearly
independent (their values at the points $\pm e_j$ distinguish them), so the eigenspace has dimension at least $n + 1$. (In fact it is exactly $n + 1$, by the theory of spherical harmonics.)
\end{solution}

\begin{exercise}\label{hodge-theorem:exercise-twisted-circle}
On the circle $S^1 = \RR/2\pi\ZZ$ let $L_a$ be the trivial complex line bundle with the flat connection $\nabla = d + ia\,d\theta$, $a \in \RR$, and the standard Hermitian metric. Compute the cohomology of the twisted
de Rham complex $0 \to C^\infty(S^1, \CC) \xrightarrow{d_\nabla} \Omega^1(S^1, \CC) \to 0$ of Example~\ref{hodge-theorem:example-elliptic-complexes}(2), and check that its Euler characteristic is $0 = \chi(S^1)$.
\end{exercise}

\begin{solution}
In Fourier series, $d_\nabla(\sum_kc_ke^{ik\theta}) = \sum_ki(k + a)c_ke^{ik\theta}\,d\theta$. If $a \notin \ZZ$, every factor $i(k + a)$ is nonzero, so $d_\nabla$ is injective, and it is surjective onto $\Omega^1$: for smooth
$g\,d\theta = \sum_kb_ke^{ik\theta}\,d\theta$ the coefficients $b_k/(i(k + a))$ decay rapidly, since $|k + a| \geq \operatorname{dist}(a, \ZZ) > 0$ grows like $|k|$. So $H^0 = H^1 = 0$. If $a = -k_0 \in \ZZ$, then
$\ker d_\nabla = \CC e^{ik_0\theta}$ and the image consists of the forms with $b_{k_0} = 0$, so $H^0 \cong H^1 \cong \CC$. In both cases $\dim H^0 - \dim H^1 = 0$. (For $a \in \ZZ$ the connection is gauge equivalent to $d$
via multiplication by $e^{-ia\theta}$, which explains why the cohomology is that of the trivial local system; for $a \notin \ZZ$ the holonomy $e^{-2\pi ia}$ is nontrivial.)
\end{solution}

\begin{exercise}\label{hodge-theorem:exercise-odd-dimension-signature}
Let $M$ be a compact oriented Riemannian manifold of dimension $4l + 2$. Show that $*$ satisfies $** = -1$ on $\mathcal{H}^{2l+1}$, and deduce that $b_{2l+1}(M)$ is even.
\end{exercise}

\begin{solution}
On $k$-forms in dimension $n$, $** = (-1)^{k(n-k)}$ (Theorem~\ref{multilinear-algebra:theorem-hodge-star}); for $n = 4l + 2$ and $k = 2l + 1$ the exponent $(2l + 1)^2$ is odd. So $*$ is a complex structure on the
real vector space $\mathcal{H}^{2l+1}$ (it preserves harmonic forms because it commutes with $\Delta$), which therefore has even dimension; conclude with Theorem~\ref{hodge-theorem:theorem-hodge}. (This also follows
from the nondegenerate skew-symmetric intersection form of Definition~\ref{poincare-duality:definition-intersection-form}.)
\end{solution}

\begin{exercise}\label{hodge-theorem:exercise-eigenvalue-one-forms}
Let $M$ be compact, connected and oriented. Show that the nonzero eigenvalues of $\Delta$ on functions are eigenvalues of $\Delta$ on $1$-forms, with eigenforms $df$, and that on a compact oriented surface the
positive eigenvalues on $1$-forms are exactly the positive eigenvalues on functions, each with twice the multiplicity.
\end{exercise}

\begin{solution}
If $\Delta f = \lambda f$ with $\lambda > 0$, then $\Delta df = \lambda df$ and $df \neq 0$ (otherwise $\lambda f = \delta df = 0$). On a surface, by Proposition~\ref{hodge-theorem:proposition-eigenspace-complex},
$E^1_\lambda = dE^0_\lambda \oplus \delta E^2_\lambda$ for $\lambda > 0$, with $d \colon E^0_\lambda \to E^1_\lambda$ injective and $\delta \colon E^2_\lambda \to E^1_\lambda$ injective (by the dual argument). The Hodge star maps
$E^0_\lambda$ isomorphically onto $E^2_\lambda$, as it commutes with $\Delta$. Hence $\dim E^1_\lambda = 2\dim E^0_\lambda$, and $E^1_\lambda = 0$ if $\lambda$ is not an eigenvalue on functions.
\end{solution}

\begin{exercise}\label{hodge-theorem:exercise-self-dual-cp2}
Let $M$ be a compact oriented Riemannian $4$-manifold. Show that for every $2$-form $\alpha$, $\alpha \wedge \alpha = (|\alpha^+|^2 - |\alpha^-|^2)\,\mathrm{vol}$, where $\alpha = \alpha^+ + \alpha^-$ is the decomposition into
self-dual and anti-self-dual parts. Deduce that if $b^- = 0$ (for instance for $\CC\PP^2$ with the Fubini--Study metric, Chapter~\ref{kahler}), every closed $2$-form $\alpha$ satisfies $\int_M\alpha \wedge \alpha \geq 0$.
\end{exercise}

\begin{solution}
$\alpha^+ \wedge \alpha^- = \alpha^- \wedge *\alpha^+ = \langle\alpha^-, \alpha^+\rangle\mathrm{vol} = 0$, $\alpha^+ \wedge \alpha^+ = \alpha^+ \wedge *\alpha^+ = |\alpha^+|^2\mathrm{vol}$ and
$\alpha^- \wedge \alpha^- = -\alpha^- \wedge *\alpha^- = -|\alpha^-|^2\mathrm{vol}$; expand $(\alpha^+ + \alpha^-)^{\wedge 2}$. For a closed $\alpha$ write $\alpha = h + d\beta$ with $h$ harmonic
(Theorem~\ref{hodge-theorem:theorem-decomposition}); by Stokes, $\int\alpha \wedge \alpha = \int h \wedge h$, and if $b^- = 0$ then $h = h^+$ is self-dual (the anti-self-dual part of a harmonic form is harmonic, as $*$
commutes with $\Delta$), so $\int h \wedge h = \|h\|^2 \geq 0$.
\end{solution}
